% !TeX root = ../../Book2.tex
\section{Poisson 方程}
\label{b2:sec:1503}

对 Poisson 方程，把二阶导数经分部积分转到测试函数上，就得到关于梯度的双线性等式。Lax--Milgram 定理可以在有界区域的能量空间中求解这一等式，并允许源项属于 \(H^{-1}\)。这种写法将存在性、稳定性与边界数据放在同一个框架中。

本节使用\subsecref{b2:subsec:150104}的实空间约定：\(H^{-1}(\Omega)\) 是带梯度范数的 \(H_0^1(\Omega;\mathbb R)\) 的连续对偶，\(\langle f,v\rangle\) 表示对偶配对。下面常数为一的估计与这一梯度范数约定相配；换用完整 \(H^1\) 范数时，对偶空间不变，数值常数则须作相应比较。

\subsection{弱解存在性}
\label{b2:subsec:150301}

\Needspace{6\baselineskip}
\begin{theorem}[Dirichlet 问题的弱可解性]
\label{b2:thm:15-poisson-existence}
设整数 \(d\geq1\)，\(\Omega\subset\mathbb R^d\) 为非空有界开集，空间取实数域。以梯度范数赋予 \(H_0^1(\Omega)\) 范数，令 \(H^{-1}(\Omega)=(H_0^1(\Omega))^*\)，其对偶范数为 \(\|f\|_{H^{-1}}=\sup_{v\in H_0^1,\,\|\nabla v\|_2\leq1}|\langle f,v\rangle|\)。对每个 \(f\in H^{-1}(\Omega)\)，存在唯一 \(u\in H_0^1(\Omega)\)，使
\[\int_\Omega\nabla u\cdot\nabla v\,\mathrm dx=\langle f,v\rangle
 \qquad(v\in H_0^1(\Omega)).\]
并且 \(\lVert\nabla u\rVert_2\le\lVert f\rVert_{H^{-1}}\)。
\end{theorem}
\begin{proof}
取 \(H=H_0^1(\Omega)\)，内积为 \(\langle u,v\rangle_H=\int\nabla u\cdot\nabla v\)。第一卷定理24.1.2的零边界 Poincaré 不等式适用于任意包含于有限宽条带的开集，故适用于这里的有界 \(\Omega\)，不要求边界正则或连通。它使梯度范数与完整 Sobolev 范数等价，故 \(H\) 完备。型 \(a(u,v)=\langle u,v\rangle_H\) 的有界常数与强制常数均为 \(1\)，右端按 \(H^{-1}\) 定义连续。\cref{b2:thm:15-lax-milgram} 因而给出存在、唯一性及估计。
\end{proof}

本例中 Lax--Milgram 退化为 Riesz 表示；一般系数时，同一个定理才显出更大的作用。该弱解也恰为 \(\tfrac12\lVert\nabla v\rVert_2^2-\langle f,v\rangle\) 的唯一极小元：只要将上一章的二次展开中的 \(L^2\) 配对换成对偶配对，证明逐字成立。两种构造在对称情形相遇，但弱型方法还可以继续处理非对称问题。

\ProseParagraphBreak

对 \(\Omega=(0,1)\)、\(f=\delta_s\)、\(0<s<1\)，解为
\[
 u(x)=\begin{cases}(1-s)x,&0\le x\le s,\\s(1-x),&s\le x\le1.\end{cases}
\]
它连续、分段仿射，导数在 \(s\) 处向下跳 \(1\)。两段分部积分给出 \(\int_0^1u'v'=v(s)\)，先对光滑 \(v\) 成立，再由点值的连续性推广。因此点源产生一个合法弱解，却不产生经典 \(C^2\) 解；这不是存在性定理的缺陷，而是数据本身的粗糙性。

\subsection{唯一性}
\label{b2:subsec:150302}

唯一性也可以直接从能量等式看出。若 \(u_1,u_2\) 具有同一数据，差 \(w\in H_0^1\) 满足 \(\int\nabla w\cdot\nabla v=0\)。取 \(v=w\)，得到 \(\nabla w=0\)；Poincaré 不等式推出 \(w=0\)。这个结论对不连通区域也成立，因为零边界约束在每个分支上排除了常数自由度。

\ProseParagraphBreak

内部方程则没有唯一性。例如任一调和函数都可加到 Poisson 方程的一个局部解上，而不改变源项。纯 Neumann 问题即使规定通量，仍允许每个连通分支上的加法常数。因而“取差再测试”最后得出梯度为零之后，还必须指出究竟由哪项边界条件或归一化排除了常数，不能直接省略最后一步。

\Needspace{6\baselineskip}
\begin{proposition}[弱比较原理]
\label{b2:prop:15-weak-comparison}
设整数 \(d\geq1\)，\(\Omega\subset\mathbb R^d\) 为非空有界开集，空间取实数域，\(H^{-1}(\Omega)=(H_0^1(\Omega))^*\)。设 \(f_1,f_2\in H^{-1}(\Omega)\)，\(u_1,u_2\in H_0^1(\Omega)\) 满足
\[
 \int_\Omega\nabla u_i\cdot\nabla v=\langle f_i,v\rangle
 \quad(v\in H_0^1(\Omega),\ i=1,2).
\]
若 \(\langle f_1-f_2,v\rangle\le0\) 对每个几乎处处非负的 \(v\in H_0^1(\Omega)\) 成立，则 \(u_1\le u_2\) 几乎处处。
\end{proposition}
\begin{proof}
令 \(w=u_1-u_2\)。第一卷命题22.2.3保证取正部保持实空间 \(H_0^1\)，命题21.3.5给出相应弱梯度公式。因此 \(w_+=\max\{w,0\}\in H_0^1\)，且 \(\nabla w_+=\mathbf1_{\{w>0\}}\nabla w\)。以它测试差方程，得
\[\lVert\nabla w_+\rVert_2^2=\langle f_1-f_2,w_+\rangle\le0.\]
零边界 Poincaré 不等式推出 \(w_+=0\)，即所需结论。所引截断结果的关键是先光滑逼近正部函数，在零水平集上用弱梯度为零处理折点，再利用正部在 \(H^1\) 中的强连续性传递零边界逼近；没有要求解逐点可微。
\end{proof}

这里的非负性还须在数据空间中解释。对非空有界开集 \(\Omega\)，实泛函 \(f\in H^{-1}(\Omega)\) 满足 \(f\geq0\)，是指
\[
 \langle f,v\rangle\geq0
 \quad\text{对所有几乎处处非负的 }v\in H_0^1(\Omega;\mathbb R).
\]
它是对测试函数作用的正性，不是给一般分布指定逐点值。若 \(f\in L^2(\Omega;\mathbb R)\) 且 \(f\geq0\) 几乎处处，已建立的 \(L^2\hookrightarrow H^{-1}\) 将配对写成可积的 \(\int fv\)，因而符合这个定义。不能只写 \(f\in L^1\) 就断言它属于 \(H^{-1}\)，也不能默认 \(fv\) 对全部 \(H_0^1\) 测试函数可积。将弱比较原理用于零源项与非负源项，得到非负数据的零边界弱解非负，与\chapref{b2:ch:13}经典最大值原理的符号一致。比较结论是几乎处处关系；尚未得到连续代表元时，不应改说“每个点都非负”。

\subsection{能量估计与 Galerkin 近似}
\label{b2:subsec:150303}

取解自身作测试，Cauchy--Schwarz 的对偶形式给出
\[\lVert\nabla u\rVert_2^2=\langle f,u\rangle
 \le\lVert f\rVert_{H^{-1}}\lVert\nabla u\rVert_2.\]
若梯度非零可约去一因子；若为零估计显然成立。因此解映射 \(H^{-1}\to H_0^1\) 连续。对于两组数据，同一计算给出
\[\lVert\nabla(u_1-u_2)\rVert_2\le\lVert f_1-f_2\rVert_{H^{-1}}.\]
特别地，平滑数据的近似会在能量范数中产生相应解的收敛，这为从光滑近似过渡到粗糙数据提供了稳定性依据。

\ProseParagraphBreak

若 \(f\in L^2\)，进一步有 \(\lVert\nabla u\rVert_2\le C_P\lVert f\rVert_2\)、\(\lVert u\rVert_2\le C_P^2\lVert f\rVert_2\)。注意区域放大时 \(C_P\) 也会变化：在 \(r\Omega\) 上换元表明 \(C_P\) 按长度尺度增长。这与二阶方程解相对于源项多出两个长度因子的事实一致，估计中的几何依赖不能无故删除。

\Needspace{6\baselineskip}
\begin{proposition}[有限维近似的最佳误差]
\label{b2:prop:15-galerkin-poisson}
设整数 \(d\geq1\)，\(\Omega\subset\mathbb R^d\) 为非空有界开集，空间取实数域，\(H^{-1}(\Omega)=(H_0^1(\Omega))^*\)。给定 \(f\in H^{-1}(\Omega)\) 与有限维线性子空间 \(V\subset H_0^1(\Omega)\)，设 \(u\in H_0^1(\Omega)\) 满足 \(\int_\Omega\nabla u\cdot\nabla v=\langle f,v\rangle\) 对全部 \(v\in H_0^1(\Omega)\) 成立。则存在唯一 \(u_V\in V\) 满足 \(\int_\Omega\nabla u_V\cdot\nabla v=\langle f,v\rangle\) 对所有 \(v\in V\) 成立，而且
\[\lVert\nabla(u-u_V)\rVert_2=\inf_{v\in V}\lVert\nabla(u-v)\rVert_2.\]
\end{proposition}
\begin{proof}
限制到 \(V\) 后型仍有界、强制，故有限维方程唯一可解。两式相减得 \(\int\nabla(u-u_V)\cdot\nabla v=0\) 对所有 \(v\in V\) 成立。对任意 \(v\in V\)，分解 \(u-v=(u-u_V)+(u_V-v)\)，两项在梯度内积下正交，故误差平方等于两项范数平方之和。最小值恰在 \(v=u_V\) 达到。
\end{proof}

近似解的误差受该有限维空间中的最佳逼近误差控制；因此，选取怎样的近似空间至关重要。若一列子空间递增且其并在 \(H_0^1\) 中稠密，最佳误差便趋零。递增性可改成对每个固定函数都能沿整个序列逼近的条件，但只有并集稠密不够，例如在良好空间之间反复插入零空间便会破坏收敛。存在性、离散方程的可解性与近似空间的逼近能力仍是三个可以分别核验的事实。

\ProseParagraphBreak

Poisson 型中，误差与近似空间真正正交，因而得到常数为一的最佳逼近。一般非对称型仍有从两个弱方程相减得到的 Galerkin 正交性，但它不能随意交换两个变量。下面的\term[cea-yinli]{Céa 引理}[Céa's lemma]用有界性和强制性代替勾股恒等式，给出准最佳逼近。

\Needspace{6\baselineskip}
\begin{proposition}[Céa 引理]\label{b2:prop:15-cea}
设 \(H\) 为实 Hilbert 空间，\(a:H\times H\to\mathbb R\) 为双线性型，存在 \(0\leq M<\infty\)、\(\alpha>0\)，使
\[
 |a(x,y)|\leq M\|x\|_H\cdot\|y\|_H,\qquad
 a(x,x)\geq\alpha\|x\|_H^2\quad(x,y\in H).
\]
给定 \(\ell\in H^*\) 及闭线性子空间 \(V\subset H\)，令 \(u\in H\)、\(u_V\in V\) 分别满足
\[
 a(u,y)=\ell(y)\quad(y\in H),\qquad
 a(u_V,v)=\ell(v)\quad(v\in V).
\]
则两解都存在且唯一，并且
\[
 \|u-u_V\|_H\leq\frac M\alpha\inf_{v\in V}\|u-v\|_H.
\]
特别地，每个有限维线性子空间均满足这里的闭性要求。
\end{proposition}
\begin{proof}
闭子空间 \(V\) 以继承的内积仍为 Hilbert 空间，限制后的型保留有界性与同一个强制常数，\(\ell|_V\) 连续。分别在 \(H\) 与 \(V\) 上使用\cref{b2:thm:15-lax-milgram}，得到两解。令 \(e=u-u_V\)，相减得到 \(a(e,v)=0\) 对所有 \(v\in V\) 成立。对任意 \(v\in V\)，因 \(v-u_V\in V\)，有
\[
 \alpha\|e\|_H^2\leq a(e,e)
 =a(e,u-v)+a(e,v-u_V)=a(e,u-v)
 \leq M\|e\|_H\cdot\|u-v\|_H.
\]
若 \(e=0\)，结论成立；否则约去 \(\|e\|_H\)，再对 \(v\in V\) 取下确界。误差关系 \(a(e,v)=0\) 对第二变量中的 \(v\in V\) 成立，这就足以得到同样适用于非对称型的估计。
\end{proof}

因此，强制型的离散解能否沿一列空间收敛，仍由这些空间能否逼近真解决定；有界性与强制性只增加了准最佳常数 \(M/\alpha\)。下面的分段仿射计算采用 Poisson 对称型，可以继续使用前一命题中常数为一的精确最佳逼近。

\ProseParagraphBreak

在区间 \((0,1)\) 上，这个结论可以算得很具体。取整数 \(N\geq2\)、均匀节点 \(x_j=jh\)（\(0\leq j\leq N\)）、\(h=1/N\)，令 \(V_h\) 为节点间分段仿射、两端为零的连续函数空间。基函数 \(\phi_j\)（\(1\leq j\leq N-1\)）在 \(x_j\) 取值一，在其余节点取值零。每个基函数只有相邻两段上的斜率非零，直接积分得
\[\int_0^1\phi_j'^2=\frac2h,\qquad
 \int_0^1\phi_j'\phi_{j+1}'=-\frac1h,\qquad
 \int_0^1\phi_j'\phi_k'=0\quad(|j-k|>1).\]
因此弱解的有限维近似由三对角线性方程组确定。矩阵的正定性并不需要另作行列式计算：任意系数向量对应的二次型就是分段仿射函数的梯度平方积分，零值只可能对应零函数。

\ProseParagraphBreak

例如取 \(f=1\)，精确弱解为 \(u(x)=x(1-x)/2\)。它在各小区间的导数平均等于两端函数值之差除以 \(h\)，而任意 \(v_h\in V_h\) 的导数在该段为常数。故 \(u\) 的节点线性插值 \(I_hu\) 满足 \(\int(u'- (I_hu)')v_h'=0\)，由唯一性即知 \(u_h=I_hu\)。在长度 \(h\) 的一段上，导数误差是以中点为零的线性函数，其平方积分为 \(h^3/12\)；函数误差则为 \((x-x_j)(x_{j+1}-x)/2\)。逐段积分并相加得到
\[\lVert u'-u_h'\rVert_2=\frac h{\sqrt{12}},\qquad
 \lVert u-u_h\rVert_2=\frac{h^2}{\sqrt{120}}.\]
这里第二个收敛速度比第一个快，来自二次函数的具体结构，不能仅从抽象存在定理读出。一般数据下的误差阶还需要解的正则性和近似空间的逼近估计；下一章提供其中的正则性工具。

\ProseParagraphBreak

弱解的误差也可以用残差直接测量。设整数 \(d\geq1\)、\(\Omega\subset\mathbb R^d\) 为非空有界开集，实值 \(f\in H^{-1}(\Omega)\)，\(u\in H_0^1(\Omega)\) 满足 \(\int_\Omega\nabla u\cdot\nabla v=\langle f,v\rangle\) 对所有实值 \(v\in H_0^1(\Omega)\) 成立。对任意实值 \(w\in H_0^1(\Omega)\) 定义 \(r=f+\Delta w\in H^{-1}\)，则 \(\langle r,v\rangle=\int\nabla(u-w)\cdot\nabla v\)。以 \(H_0^1\) 的梯度范数取对偶范数，对单位梯度范数的测试函数取上确界，Cauchy--Schwarz 给出上界，而在误差非零时取 \(v=(u-w)/\lVert\nabla(u-w)\rVert_2\) 达到该界。因此 \(\lVert r\rVert_{H^{-1}}=\lVert\nabla(u-w)\rVert_2\)。对分段仿射近似，残差可能含节点通量跳跃，不能只计算每段内为零的二阶导数；负范数恰好保留这些分布贡献。

\subsection{非齐次 Dirichlet 数据与 Neumann 问题}
\label{b2:subsec:150304}

取整数 \(d\geq1\)、非空有界开集 \(\Omega\subset\mathbb R^d\)、实值 \(f\in H^{-1}(\Omega)\) 与 \(g\in H^1(\Omega;\mathbb R)\)。令 \(u=g+w\)，其中 \(w\in H_0^1(\Omega;\mathbb R)\)。方程等价于
\[\int_\Omega\nabla w\cdot\nabla v
  =\langle f,v\rangle-\int_\Omega\nabla g\cdot\nabla v\qquad(v\in H_0^1).\]
右端属于 \(H^{-1}\)，范数不超过 \(\lVert f\rVert_{H^{-1}}+\lVert\nabla g\rVert_2\)。因此齐次理论给出唯一 \(w\)，从而给出唯一满足 \(u-g\in H_0^1\) 的解；并有
\[\lVert\nabla w\rVert_2\le\lVert f\rVert_{H^{-1}}+\lVert\nabla g\rVert_2.\]
完整 \(H^1\) 估计还要加上 \(g\) 的 \(L^2\) 范数。

\ProseParagraphBreak

若另一个延拓 \(\widetilde g\) 与 \(g\) 相差 \(H_0^1\) 函数，两者给出的仿射可行集完全相同；由唯一性得到同一个 \(u\)。所以解依赖边界类，而不依赖人为选择的延拓。若要比较两组边界数据，先选择它们差的一个延拓，再把上面的稳定性估计应用于差方程；优化延拓的范数便自然引出商空间形式的迹范数。

\ProseParagraphBreak

例如在 \(\mathbb R^d\) 的单位球上给定仿射边界数据 \(g(x)=a\cdot x+b\)，其中 \(a\in\mathbb R^d\)、\(b\in\mathbb R\)，且 \(f=0\)。由于 \(\int a\cdot\nabla v=0\) 对 \(v\in H_0^1\) 成立，齐次修正为零，故 \(u=g\)。对一般边界数据，弱方法也能构造能量解；在球上还可借\chapref{b2:ch:13}的表示公式识别这个解。不同构造的解一旦处于同一能量边界类，就由唯一性自动一致。


\ProseParagraphBreak

\noindent\textbf{Neumann 问题。}区域保持固定而边界值自由时，应换一个 Hilbert 空间，不能沿用零 Dirichlet 条件排除常数的唯一性论证。

\Needspace{6\baselineskip}
\begin{proposition}[Neumann 数据与归一化]
\label{b2:prop:15-neumann-normalized}
设整数 \(d\geq1\)，\(\Omega\subset\mathbb R^d\) 为非空有界连通 Lipschitz 区域，\(f\in L^2(\Omega;\mathbb R)\)、\(g_N\in L^2(\partial\Omega;\mathbb R)\)，边界积分取 \((d-1)\) 维面积测度（\(d=1\) 时为端点计数测度），\(\operatorname{Tr}\) 为 Sobolev 迹，求解类为 \(u\in H^1(\Omega;\mathbb R)\)。方程
\[\int_\Omega\nabla u\cdot\nabla v
 =\int_\Omega fv+\int_{\partial\Omega}g_N\operatorname{Tr}v
 \quad(v\in H^1(\Omega))\]
可解当且仅当 \(\int_\Omega f+\int_{\partial\Omega}g_N=0\)。指定 \(\int_\Omega u=0\) 后解唯一，且
\[\|u\|_{H^1}\le C_\Omega(\|f\|_2+\|g_N\|_{L^2(\partial\Omega)}).\]
\end{proposition}
\begin{proof}
测试 \(v=1\) 给出必要条件。令 \(H=\{v\in H^1(\Omega;\mathbb R):\int_\Omega v=0\}\)。积分泛函因 \(\Omega\) 测度有限而连续，故 \(H\) 是闭子空间；第一卷定理24.1.3的 Poincaré--Wirtinger 不等式适用于这里的有界连通 Lipschitz 区域，使 \(\|\nabla v\|_2\) 成为 \(H\) 上等价的完备范数。右端有界性来自
\[
 |\ell(v)|\le\|f\|_2\|v\|_2+
 \|g_N\|_{L^2(\partial\Omega)}\|\operatorname{Tr}v\|_{L^2(\partial\Omega)}
 \le C_\Omega(\|f\|_2+\|g_N\|_2)\|\nabla v\|_2,
\]
这里第一卷定理23.1.3给出有界 Lipschitz 区域上的 \(H^1\to L^2(\partial\Omega)\) 迹估计，刚才的零均值不等式再将 \(H^1\) 范数换成梯度范数。Lax--Milgram 在 \(H\) 上给出唯一解及估计。对任意 \(v\in H^1(\Omega)\)，分解 \(v=v_0+v_\Omega\)，其中 \(v_\Omega=|\Omega|^{-1}\int_\Omega v\)、\(v_0\in H\)。左端常数项为零，右端常数项因相容条件为零，所以受限测试解确实解全部测试方程。两解之差测试自身，得到梯度为零；连通性说明差为一个常数，归一化排除它。
\end{proof}

在 \((0,1)\) 上取 \(f=1\)、左端外向通量 \(g_N(0)=-1\)、右端通量 \(g_N(1)=0\)。总数据为 \(1-1+0=0\)。积分 \(-u''=1\) 得 \(u'=-x+C\)；左端条件 \(-u'(0)=-1\) 给 \(C=1\)，右端也满足。减去均值后 \(u=x-x^2/2-1/3\)。这个例子同时检验了相容条件中的加号和左端外法向的负号。若区域有多个连通分支，必须在每个分支分别相容、分别归一化，单个全局零均值无法排除所有常数自由度。
